% =====================================================================
%  2bat-ud6-12.tex  —  SESSIÓ 12: Detectar lectures enganyoses
%  (sense preàmbul: s'inclou des de main.tex)
% =====================================================================
\horatitol{Sessió 12 — Detectar lectures enganyoses}%
{Activitat 4 (fase 2): contrastar titulars de premsa amb les dades reals dels estudis i construir un catàleg de paranys a evitar.}

\subsection{Idea clau}

Completem l'Activitat 4 centrant-nos en com els resultats d'un estudi es poden
\textbf{distorsionar} en comunicar-los. A les \textbf{pissarres}, en grups, contrastareu
titulars de premsa amb les dades reals de l'estudi i decidireu si el titular diu realment
el que diuen les dades.

\begin{idea}
\textbf{Els paranys més habituals en comunicar dades}
\begin{itemize}[leftmargin=1.4em, itemsep=2pt]
  \item \textbf{Correlació presentada com a causa:} dues coses van juntes i es conclou que
        una \emph{causa} l'altra.
  \item \textbf{Ignorar el marge d'error:} presentar una estimació com una certesa exacta.
  \item \textbf{Generalitzar des d'una mostra esbiaixada:} estendre a tothom el resultat
        d'un grup no representatiu.
  \item \textbf{Inflar una diferència no significativa:} destacar com a rellevant una
        diferència que pot ser atzar.
  \item \textbf{Triar la comparació que convé:} escollir la dada o el període que dona el
        titular més cridaner.
\end{itemize}
\textbf{La pregunta clau:} el titular diu \emph{realment} el que diuen les dades?
\end{idea}

\subsection{Exemples}

\begin{exemple}[correlació convertida en causa]
Titular: «\emph{Els joves que fan activitats al centre cívic treuen millors notes}». Dada
real: s'ha observat que els que hi van tendeixen a tenir millors notes. \textbf{Parany:}
d'una \textbf{correlació} (van juntes) es conclou una \textbf{causa} (el centre millora les
notes). Podria ser al revés (els que ja estudien més també van al centre) o haver-hi una
tercera causa (famílies implicades). El titular sobreinterpreta.
\end{exemple}

\begin{exemple}[ignorar el marge d'error]
Titular: «\emph{Puja el suport al servei: del $48\%$ al $51\%$}». Dada real: totes dues
enquestes tenen un marge de $\pm 3\%$. \textbf{Parany:} els intervals ($[45\%,51\%]$ i
$[48\%,54\%]$) se superposen, així que la «pujada» de $3$ punts pot ser atzar. Presentar-la
com una tendència real ignora el marge d'error.
\end{exemple}

\subsection{Exercicis resolts}

\begin{exresolt}[detectar el parany]
Titular: «\emph{Un estudi demostra que la població prefereix el servei A}», a partir d'una
enquesta clicable en una web. Quin parany hi ha?
\solucio
Doble parany: (1) \textbf{generalitzar des d'una mostra esbiaixada} —una enquesta clicable
recull només qui s'autoselecciona, no la població; (2) el verb «\textbf{demostra}» presenta
com a certesa el que és, com a molt, una estimació d'una mostra poc fiable. El titular no es
pot sostenir amb aquestes dades.
\resposta{Mostra autoseleccionada + «demostra» (certesa infundada).}
\end{exresolt}

\begin{exresolt}[reescriure el titular amb honestedat]
Reescriu de manera honesta el titular «\emph{Puja el suport del $48\%$ al $51\%$}» (marges
de $\pm 3\%$).
\solucio
Una versió honesta seria: «\emph{El suport es manté estable al voltant del $50\%$: la
variació observada està dins del marge d'error}». Així s'evita presentar com a tendència una
diferència que pot ser fruit de l'atzar.
\resposta{Reconèixer que la variació cau dins del marge i que el suport es manté estable.}
\end{exresolt}

\subsection{Exercicis per fer}

\noindent\textit{Treballeu en grup a la pissarra: per a cada titular, digueu quin parany hi
ha i com el corregiríeu.}

\begin{exercicis}
\begin{exlist}
  \item Titular: «\emph{Els municipis amb més biblioteques tenen menys atur}». Dada: s'ha
        observat una correlació. Quin parany hi ha?
  \item Titular: «\emph{El $70\%$ de la població dona suport a la mesura}», a partir d'una
        enquesta feta a la sortida d'un acte a favor de la mesura. Quin biaix?
  \item Titular: «\emph{Rècord de satisfacció: puja del $82\%$ al $83\%$}», amb marges de
        $\pm 2\%$. La pujada és real o pot ser atzar?
  \item Titular: «\emph{Un estudi demostra que el programa cura l'ansietat}», sense grup de
        comparació. Quin problema té la paraula «demostra»?
  \item \textbf{(Reescriptura)} Tria un dels titulars anteriors i reescriu-lo de manera
        fidel a les dades.
  \item \textbf{(Catàleg de paranys)} Escriu la teva llista dels paranys més habituals a
        evitar quan es llegeixen o es comuniquen dades.
\end{exlist}
\end{exercicis}

% ---------------------------------------------------------------------
\fullsolucions{Sessió 12}
\begin{solucions}
\begin{sollist}
  \item Correlació presentada com a causa: que dues coses vagin juntes (més biblioteques,
        menys atur) no vol dir que una causi l'altra; pot haver-hi una tercera causa (per
        exemple, municipis més rics tenen més biblioteques \emph{i} menys atur).
  \item Biaix de selecció: s'enquesta qui assisteix a un acte \emph{a favor}, un públic ja
        favorable. El $70\%$ no és generalitzable a tota la població.
  \item Pot ser atzar: amb marges de $\pm 2\%$, els intervals ($[80\%,84\%]$ i
        $[81\%,85\%]$) se superposen; $1$ punt de «pujada» no és una tendència fiable.
  \item «Demostra» presenta com a certesa causal el que, sense grup de comparació, no es pot
        establir: potser l'ansietat hauria baixat igualment. Caldria dir, com a molt, que
        «els participants milloren», sense afirmar causa provada.
  \item Resposta oberta. Ha de reescriure el titular expressant la incertesa i evitant el
        parany (per exemple, substituir «demostra» per «suggereix», o reconèixer que la
        variació cau dins del marge).
  \item Resposta oberta. Ha d'incloure: confondre correlació amb causa, ignorar el marge
        d'error, generalitzar des d'una mostra esbiaixada, inflar diferències no
        significatives i triar la comparació que convé.
\end{sollist}
\end{solucions}
